\section{Introduction}

When an Indian examination board declares results, lakhs of students try to
open the result website at the same time~\cite{ht2021ssc}. The release time is
announced in advance, almost every request asks for one roll number, and the
results do not change while students read them. Few web workloads are
simpler, yet result websites keep failing. The chairman of JEE Advanced 2016
described the portal receiving about \num{10000} requests per second against
a server that could hold \num{2000}~\cite{pagalguy2016jee}. On Maharashtra SSC
result day in 2021, schools reported that at least half of their students
could not see results until evening~\cite{ht2021ssc}, and the state's
education minister apologised~\cite{careerindia2021ssc}. Kenya's examination
council reported up to \num{300000} hits per second~\cite{knec2025x}, and in
2025 WAEC in Nigeria withdrew its result checker over \emph{wrong}
results~\cite{icir2025waec}.

A board that wants to avoid this has to decide how to build its website, and
the options are well known: caching or pre-generating
pages~\cite{challenger1999scalable,labrinidis2000webview}, protecting caches
against stampedes~\cite{vattani2015xfetch,nishtala2013memcache}, admission
control~\cite{welsh2001seda,elnikety2004gatekeeper,chen2005news}, and
``virtual waiting rooms'' that queue visitors during surges, sold as
products~\cite{cloudflare2026wr,aws2022vwr} and studied in WebQ, which tells
users of an overloaded server how long to wait and checks with a token that
they did~\cite{doshi2015webq}. What is missing is evidence about which option
works for this workload. Flash-crowd studies characterise crowds and absorb
them with caching~\cite{ari2003flash,jung2002flash,wendell2011going}, but
these crowds are often unforeseen, while a result release is announced. The closest recent studies compare a database-only design with a
Redis cache for university systems under a fixed number of users, one with a
real database~\cite{putranto2026redis} and one with simulated
delays~\cite{tofan2026cacheaside}; neither tests pre-computed results or
waiting rooms, models students who refresh, or checks that the answers are
correct. To our knowledge, no study has compared waiting rooms with other
designs under an announced surge of students who refresh.

This paper gives boards measured evidence for four decisions: which
architecture to use (Q1), whether to put a waiting room in front of it (Q2),
how to keep results correct when they are corrected or released early (Q3),
and how caches and refreshing students change the load (Q4). We contribute
ResultBench, an open-source testbed that runs eleven designs under one
resource budget, replays a result-day surge of students who refresh, and
checks every answer against ground truth; a controlled comparison of the
designs; and a decision guide that turns the measurements into advice
(Table~\ref{tab:guide}).
